\documentclass[a4paper]{amsart}
\usepackage[a4paper, left=2cm, right=2cm, top=2cm, bottom=2cm]{geometry}

\usepackage{graphicx, float, enumerate, amsmath, amsfonts, amssymb, amsthm, xcolor, url, hyperref, multicol, accents,footnote}
\usepackage[french]{babel}

% \usepackage{enumerate, amsmath, amsfonts, amssymb, amsthm, wasysym, graphics, graphicx, xcolor, url, hyperref, hypcap, a4wide, stmaryrd, shuffle, xargs, multicol, overpic, pdflscape, multirow, hvfloat, minibox, accents, etoolbox, dsfont}

\makeatletter
% we use \prefix@<level> only if it is defined
\renewcommand{\@seccntformat}[1]{%
  \ifcsname prefix@#1\endcsname
    \csname prefix@#1\endcsname
  \else
    \csname the#1\endcsname\quad
  \fi}
% define \prefix@section
\newcommand\prefix@section{\thesection}
\makeatother

\renewcommand{\thesection}{}
\renewcommand{\thesubsection}{\arabic{subsection}}

\newcommand{\blue}{blue!65!white}
\newcommand{\mydef}[1]{\textcolor{\blue}{{\bf #1}}} % blue bold definition
% % \newcommand{\mydef}[1]{{\bf #1}} % bold definition


\title{Treillis join-distributifs}


\addto\captionsfrench{% Replace "english" with the language you use
  \renewcommand{\contentsname}%
    {Contenu du document}%
}

\newtheorem{lem}{Lemme}
\newtheorem{theo}{Théorème}
\newtheorem{defi}{Définition}
\newtheorem{prop}{Proposition}
\newtheorem{conj}{Conjecture}
\newtheorem{rmk}{Remarque}

\newcommand{\dps}{\displaystyle}
\renewcommand{\L}{\mathbf L}
\renewcommand{\P}{\mathcal P}
\newcommand{\InfPP}{\mathcal Y_{\infty}}
\newcommand{\ChainP}{\mathcal C}
\newcommand{\LS}{\downarrow \!}
\newcommand{\US}{\uparrow \!}
\newcommand{\Pos}{\mathcal P}
\newcommand{\Anti}{\operatorname{Anti}}
\newcommand{\Dyck}{\mathbf{Dyck}}
\newcommand{\Motz}{\mathbf{Motz}}
\newcommand{\Schr}{\mathbf{Schr}}
\newcommand{\MD}{\mathbf{MD}}
\newcommand{\CS}{\operatorname{CS}}
\newcommand{\stand}{\operatorname{St}}
\newcommand{\Ss}{\mathfrak{S}}
\newcommand{\N}{\mathbb{N}}
\newcommand{\Z}{\mathbb{Z}}
\newcommand{\Irr}{\operatorname{Irr}}
\newcommand{\JIrr}{\mathcal J}
\newcommand{\MIrr}{\mathcal M}
\newcommand{\IdL}{\mathbf{Low}}
\newcommand{\Dis}{\operatorname{Dis}}
\newcommand{\Con}{\operatorname{Con}}
\newcommand{\AD}{\operatorname{AscDes}}
\newcommand{\Dec}{\operatorname{Dec}}

\newcommand{\labels}{\Lambda}


\begin{document}
\maketitle

\section{Introduction}

Citer \cite{AFN22} \cite{Czed14} \cite{Arm09}

\begin{theo}
   Soit $\L$ un treillis fini. Les assertions suivantes sont équivalentes :
   \begin{enumerate}
   \item \label{item:1} Tout intervalle atomique de $\L$ est booléen ;
   \item \label{item:2} Tout élément de $\L$ s'écrit de manière unique comme le \emph{meet} d'une antichaîne d'éléments \emph{meet}-irréductibles ;
   \item \label{item:3} $\L$ est semimodulaire supérieur et \emph{meet}-semidistributif ;
   \item \label{item:4} toutes les chaînes maximales de $\L$ ont pour longueur $|\MIrr(\L)|$ ;
   \item \label{item:5} $\L$ est un \emph{join}-sous-semitreillis d'un treillis distributif fini ;
   \item \label{item:6} $\L$ est isomorphe au treillis des ouverts d'une clôture d'antiéchange (ou antimatroïde, ou géométrie convexe).
   \end{enumerate}
\end{theo}

Un treillis qui vérifie ces propriétés est dit \emph{join}-distributif. Un treillis dont le dual est \emph{join}-distributif est dit \emph{meet}-distributif.

Un treillis \emph{join}-distributif étant \emph{meet}-semidistributif, $\L$ est gradué de hauteur $|\MIrr(\L)|$. Il possède une application $\kappa:\JIrr(\Pos)\to \MIrr(\Pos)$ définie par $\kappa(j)=\max\{x\in \L:x\wedge j=j_*\}$.

L'assertion (\ref{item:6}) est analogue au résultat suivant :
\begin{rmk}
Un treillis fini est géométrique si et seulement si il est isomorphe au treillis des ensembles fermés d'une clôture d'échange.
\end{rmk}

\section{Premiers résultats}

\begin{prop}
Soit $\Pos$ un poset fini, $n\in \N$ et $\lambda :\Pos \to \{1,...,n\}$ un étiquetage de $\Pos$.
$\IdL(\Pos,\lambda):=\{\lambda(X):X\in \IdL(\Pos)\}$ ordonné par l'inclusion est un treillis \emph{join}-distributif.
\end{prop}
\begin{proof}
  $\IdL(\Pos,\lambda)$ est stable par union donc c'est un \emph{join}-semitreillis, et il possède un plus grand élément $\lambda(\Pos)$, c'est donc un treillis~\cite[Proposition 3.3.1]{Stan12}. Les relations de couvertures de $\IdL(\Pos,\lambda)$ sont de la forme $\lambda(X)\lessdot \lambda(X)\cup\{k\}$, et comme le \emph{join} dans $\IdL(\Pos,\lambda)$ est l'union d'ensembles on en déduit que ses intervalles atomiques sont tous booléens.
  %Soit $X\in \IdL(\Pos)$ et $k_1,\dots,k_r$ tels que $\lambda(X)\cup\{k_1\}, \dots \lambda(X)\cup\{k_r\}$ couvrent $\lambda(X)$ dans $\IdL(\Pos,\lambda)$.
\end{proof}

L'application $\lambda :\IdL(\Pos):=\IdL(\Pos, \lambda)$ est un morphisme de \emph{join}-semitreillis : $\lambda(X)\cup\lambda(Y)=\lambda (X\cup Y)$. On en déduit le résultat suivant :
\begin{prop}
   Pour tout $j\in \JIrr(\IdL(\Pos, \lambda))$, les éléments minimaux de $\lambda^{-1}(j)$ appartiennent à $\JIrr(\IdL(\Pos))$, \emph{i.e.} ce sont des idéaux principaux de $\Pos$.
   \label{prop:2}
\end{prop}
\begin{proof}
   Supposons que $x$ soit minimal dans $\lambda^{-1}(j)$ sans être \emph{join}-irréductible. Il existe donc $y$ et $z$ couverts par $x$ dans $\IdL(\Pos)$ tels que $\lambda(y)$ et $\lambda(z)$ sont différents de $j$. Or $j=\lambda(x)= \lambda(y\cup z)=\lambda(y)\cup \lambda(z)$, et comme $j$ est \emph{join}-irréductible, $\lambda(y)$ ou $\lambda(z)$ doit être égal à $j$.
\end{proof}


\begin{prop}
   Soit $\L$ un treillis \emph{meet}-semidistributif fini. Pour tout $x\in \L$, on définit $\labels(x):=\{\kappa(j):j\in \JIrr(\L), j\le x\}$. $\labels(x)$ caractérise $x$, $x$ étant l'unique élément minimal de $\L\setminus \bigcup_{m\in \labels(x)} [\le m]$.
\end{prop}

\begin{prop}
  Soit $\L$ un treillis \emph{join}-distributif fini. $\L$ est isomorphe à $\IdL(\JIrr(\L),\kappa)$.
\end{prop}

\begin{lem}
   Soit $x<y$ dans $\Pos$ tel que $\lambda(x)=\lambda(y)$. Il existe $z\in \Pos$ tel que $\lambda([\le z])=\lambda([\le y])$, ou alors $\lambda([\le y])$ n'est pas \emph{join}-irréductible dans $\IdL(\Pos,\lambda)$.
\end{lem}
\begin{proof}
   Si $x<y$ dans $\Pos$ avec $\lambda(x)=\lambda(y)$, alors $\lambda([< y])=\lambda([\le y])$. D'après la Proposition \ref{prop:2}, il existe un idéal principal $[\le z]$ strictement inclus dans $[\le y]$ tel que $\lambda([\le z])=\lambda([\le y])$, ou alors $\lambda([\le y])$ ne serait pas \emph{join}-irréductible.
\end{proof}

Tout treillis \emph{join}-distributif $\L$ peut s'obtenir à partir d'un poset étiqueté avec autant d'éléments qu'il y a de \emph{join}-irréductibles dans $\L$. À quel condition un poset étiqueté $(\Pos,\lambda)$ est de cette forme ? Autrement dit, on veut que les $\lambda([\le x])$ soit tous distincts et \emph{join}-irréductibles dans $\IdL(\Pos, \lambda)$.


\subsection{Complétion}

\begin{prop}
   Soit $\Pos$ un poset fini. $\MD(P)$ est \emph{join}-distributif si et seulement si $(\min \{x:x\nleq m\})_{m\in \MIrr(P)}$ est une partition de $\JIrr(\Pos)$.
\end{prop}

\section{Treillis \emph{join}-distributifs semi-infinis}

Treillis des multipartitions (Izumi Watanabe)

Étendre les résultats de représentation, complétion des treillis join-distributifs aux treillis semi-infinis.

Voir "An extension of Birkhoff's representation theorem to locally-finite distributive lattices" \cite{Worl26}

\bibliography{lattices}{}
\bibliographystyle{plain}

\end{document}